\section{Implementasi RSA dalam C}

Implementasi RSA yang aman memerlukan perhatian pada side-channel dan operasi constant-time \cite{ferguson2010}. Untuk pembelajaran gunakan bilangan kecil; untuk produksi gunakan library (OpenSSL, GMP).

\subsection{Implementasi Edukatif dengan Bilangan Kecil}

Implementasi sederhana berguna untuk memahami konsep dasar RSA:
\begin{lstlisting}[style=cstyle, caption={RSA Sederhana dengan Bilangan Kecil (C)}]
#include <stdio.h>

long long mod_pow(long long base, long long exp, long long mod) {
    long long result = 1;
    base = base % mod;
    while (exp > 0) {
        if (exp % 2 == 1) result = (result * base) % mod;
        exp /= 2;
        base = (base * base) % mod;
    }
    return result;
}

int main() {
    long long p = 61, q = 53;
    long long n = p * q;
    long long phi = (p-1) * (q-1);
    long long e = 17;
    long long d = 2753;  /* e^(-1) mod phi, precomputed */
    
    long long m = 65;  /* plaintext */
    long long c = mod_pow(m, e, n);
    long long m2 = mod_pow(c, d, n);
    
    printf("Plaintext: %lld\n", m);
    printf("Ciphertext: %lld\n", c);
    printf("Decrypted: %lld\n", m2);
    return 0;
}
\end{lstlisting}

Kode di atas hanya untuk pembelajaran; tidak memiliki padding aman atau constant-time operations.

\subsection{Implementasi Produksi dengan OpenSSL}

Untuk RSA nyata, gunakan API OpenSSL berbasis \code{EVP\_PKEY} dan padding OAEP/PSS \cite{rfc8017}.

\begin{lstlisting}[style=cstyle, caption={Cuplikan EVP untuk RSA-OAEP (ringkas)}]
EVP_PKEY_CTX *ctx = EVP_PKEY_CTX_new(pkey, NULL);
EVP_PKEY_encrypt_init(ctx);
EVP_PKEY_CTX_set_rsa_padding(ctx, RSA_PKCS1_OAEP_PADDING);
EVP_PKEY_CTX_set_rsa_oaep_md(ctx, EVP_sha256());
EVP_PKEY_encrypt(ctx, out, &outlen, in, inlen);
\end{lstlisting}

\begin{lstlisting}[style=cstyle, caption={Cuplikan EVP untuk RSA-PSS (ringkas)}]
EVP_PKEY_CTX *ctx = EVP_PKEY_CTX_new(pkey, NULL);
EVP_PKEY_sign_init(ctx);
EVP_PKEY_CTX_set_rsa_padding(ctx, RSA_PKCS1_PSS_PADDING);
EVP_PKEY_CTX_set_signature_md(ctx, EVP_sha256());
EVP_PKEY_sign(ctx, sig, &siglen, msg, msglen);
\end{lstlisting}


\subsection{Keamanan Implementasi dan Side-Channel Attacks}

Implementasi RSA yang aman harus memperhatikan berbagai aspek keamanan:
\begin{itemize}
  \item \textbf{Random Number Generation}: Gunakan RNG kriptografis yang kuat untuk kunci dan salt.
  \item \textbf{Constant-Time Operations}: Implementasi operasi modular yang waktu eksekusinya konstan.
  \item \textbf{Blinding}: Gunakan cryptographic blinding untuk melindungi dari timing attacks.
  \item \textbf{Memory Security}: Hapus sensitif data dari memory setelah digunakan.
  \item \textbf{Input Validation}: Validasi semua input sebelum diproses.
\end{itemize}

Timing attacks (Kocher, 1995) dapat mendeduksi kunci dari waktu dekripsi. Gunakan constant-time dan blinding \cite{ferguson2010}.

\subsection{Cryptographic Blinding}

Cryptographic blinding melindungi dari timing attacks.

Implementasi blinding dalam RSA:
\begin{enumerate}
  \item Pilih nilai acak $r$ untuk setiap operasi dekripsi
  \item Hitung $c' = c \cdot r^e \bmod n$
  \item Lakukan dekripsi: $m' = (c')^d \bmod n$
  \item Hilangkan blinding: $m = m' \cdot r^{-1} \bmod n$
\end{enumerate}


\subsection{Optimasi Performa dengan Chinese Remainder Theorem}

Optimasi CRT memberikan percepatan hingga 4x untuk dekripsi.

\begin{lstlisting}[style=cstyle, caption={Optimasi CRT untuk RSA Dekripsi}]
// Precompute CRT parameters
long long dP = d % (p-1);
long long dQ = d % (q-1);
long long qInv = mod_inverse(q, p);

// CRT-based decryption
long long m1 = mod_pow(c, dP, p);
long long m2 = mod_pow(c, dQ, q);
long long h = ((qInv * (m1 - m2)) % p) * q;

long long m = (m2 + h) % n;
\end{lstlisting}


\subsection{Validasi dan Testing}

Implementasi RSA harus divalidasi secara menyeluruh sebelum digunakan dalam produksi:
\begin{itemize}
  \item \textbf{Known Answer Tests}: Uji dengan vektor test yang telah dipublikasikan.
  \item \textbf{Randomized Tests}: Uji dengan input acak untuk mendeteksi bug.
  \item \textbf{Side-Channel Tests}: Uji resistensi terhadap timing dan cache attacks.
  \item \textbf{Performance Tests}: Uji throughput dan latensi untuk berbagai ukuran kunci.
\end{itemize}

